% Figure from MATH 556 notes, section 25. Compile with the notes preamble.
\begin{tikzpicture}
\begin{axis}[nfig, width=6.6cm, height=6.6cm, clip=false,
  xmin=0, xmax=14, ymin=-0.56, ymax=0.15, axis x line=middle, axis y line=left,
  xtick=\empty, ytick={-0.5,-0.125,0}, yticklabels={$-\tfrac12$,$-\tfrac18$,$0$},
  xlabel={$|x|$}, ylabel={$E$},
  xlabel style={anchor=south west, font=\scriptsize, inner sep=1pt},
  ylabel style={rotate=-90, font=\scriptsize, at={(ticklabel* cs:1)}, anchor=south}]
  \fill[figR!15] (axis cs:0,0) rectangle (axis cs:14,0.15);
  \draw[figR, thick] (axis cs:0,0) -- (axis cs:14,0);
  \addplot[black!55, thick, domain=1.79:14, samples=120] {-1/x};
  \node[black!60, font=\scriptsize, anchor=west] at (axis cs:4.5,-0.33) {$V = -\dfrac{1}{|x|}$};
  % bound levels E_n = -1/(2n^2), drawn out to the classical turning point |x| = 2n^2
  \addplot[figB, thick] coordinates {(0,-0.5) (2,-0.5)};
  \addplot[figB, thick] coordinates {(0,-0.125) (8,-0.125)};
  \foreach \n in {3,...,9} {
    \edef\temp{\noexpand\addplot[figB, line width=0.55pt] coordinates {(0,{-1/(2*\n*\n)}) (14,{-1/(2*\n*\n)})};}
    \temp
  }
  \node[figB, font=\scriptsize, anchor=west] at (axis cs:2.1,-0.5) {$E_1$, dim $1$};
  \node[figB, font=\scriptsize, anchor=west] at (axis cs:8.1,-0.125) {$E_2$, dim $4$};
  \node[figB, font=\scriptsize, anchor=west] at (axis cs:14.1,-0.0556) {$E_3$, dim $9$};
  \node[figB, font=\scriptsize, anchor=west] at (axis cs:14.1,-0.017) {$E_n \uparrow 0$};
  \node[figR, font=\scriptsize, anchor=center] at (axis cs:7,0.08) {continuum $[0, \infty)$};
\end{axis}
\end{tikzpicture}
